\documentclass[11pt,a4paper]{article}
\usepackage[utf8]{inputenc}
\usepackage[T1]{fontenc}
\usepackage{lmodern}
\usepackage[brazilian]{babel}
\usepackage{amsmath,amssymb}
\usepackage{geometry}
\geometry{top=2.0cm, bottom=2.0cm, left=2.0cm, right=2.0cm}
\usepackage{enumitem}
\usepackage{fancyhdr}
\usepackage{titlesec}
\usepackage{xcolor}
\usepackage{tcolorbox}

\pagestyle{fancy}
\fancyhf{}
\fancyhead[L]{\small\textbf{FSC5101 -- Física I} $\vert$ UFSC}
\fancyhead[R]{\small\textbf{Gabarito Comentado e Didático} $\vert$ Capítulos 4 e 5}
\fancyfoot[C]{\thepage}
\renewcommand{\headrulewidth}{0.5pt}

\titleformat{\section}{\large\bfseries\color{blue!70!black}}{}{0em}{}[\titlerule]
\titleformat{\subsection}{\bfseries\color{black}}{}{0em}{}

% Caixas de destaque para dicas e respostas
\tcbset{
    boxrule=0.8pt,
    arc=2mm,
    left=3mm, right=3mm, top=2mm, bottom=2mm
}
\newtcolorbox{dica}{
    colback=yellow!10!white,
    colframe=orange!80!black,
    title=\textbf{Dica Metodológica / Pulo do Gato:}
}
\newtcolorbox{dclbox}{
    colback=blue!5!white,
    colframe=blue!60!black,
    title=\textbf{Diagrama de Corpo Livre (DCL) e Identificação de Forças:}
}

\begin{document}

\begin{center}
    {\LARGE \textbf{UNIVERSIDADE FEDERAL DE SANTA CATARINA}}\\[0.2cm]
    {\large \textbf{Departamento de Física -- FSC5101: Física I}}\\[0.3cm]
    {\Large \textbf{Gabarito Completo, Didático e Comentado (40 Questões)}}\\[0.2cm]
    {\normalsize \textbf{Fonte:} Sears \& Zemansky -- Física I: Mecânica (Young \& Freedman, 12ª Edição)}\\[0.1cm]
    {\small \textit{Capítulo 4: Leis de Newton do Movimento \quad $\vert$ \quad Capítulo 5: Aplicações das Leis de Newton}}\\[0.2cm]
    \rule{\textwidth}{0.8pt}
\end{center}

\vspace{0.2cm}

%%%%%%%%%%%%%%%%%%%%%%%%%%%%%%%%%%%%%%%%%%%%%%%%%%%%%%%%%%%%%%%%%%%%%%%%%%%%%%%
\section*{GUIA PRÁTICO: O DIAGRAMA DE CORPO LIVRE E PLANOS INCLINADOS}
%%%%%%%%%%%%%%%%%%%%%%%%%%%%%%%%%%%%%%%%%%%%%%%%%%%%%%%%%%%%%%%%%%%%%%%%%%%%%%%

Antes de iniciar as resoluções, destacamos as regras de ouro para construir o Diagrama de Corpo Livre (DCL) e trabalhar com ângulos que não começam em $0^\circ$ (especialmente em planos inclinados):

\begin{enumerate}[label=\textbf{\Roman*.}]
    \item \textbf{A Regra Básica do DCL (Isolamento do Corpo):}
    O DCL contém \textbf{apenas} as forças externas que o ambiente (outros corpos) aplica \textbf{sobre} o corpo em estudo. Nunca inclua no DCL forças que o corpo exerce sobre os outros (Terceira Lei) e nunca desenhe o vetor $m\vec{a}$ como se fosse uma força física aplicada.

    \item \textbf{Alinhamento Inteligente dos Eixos Coordenados:}
    Em superfícies inclinadas de um ângulo $\theta$ com a horizontal, alinhar os eixos com a horizontal e a vertical força você a decompor a força normal, o atrito e a aceleração. A escolha ideal é \textbf{girar os eixos}:
    \begin{itemize}
        \item \textbf{Eixo $x$ (paralelo à rampa):} aponta ao longo do movimento (para cima ou para baixo da rampa).
        \item \textbf{Eixo $y$ (perpendicular à rampa):} aponta para fora da rampa, paralelo à força normal $\vec{n}$.
    \end{itemize}
    Dessa forma, a aceleração e a força normal ficam sobre os eixos ($\sum F_y = 0$), e a única força que precisa ser decomposta é a força peso!

    \item \textbf{Onde vai o ângulo $\theta$ na decomposição do Peso?}
    O ângulo $\theta$ do plano inclinado é medido em relação à \textbf{horizontal}. Como o peso $\vec{P}$ é vertical e o eixo $y$ é perpendicular à rampa:
    \begin{itemize}
        \item A linha vertical do peso e o eixo perpendicular $y$ formam o \textbf{mesmo ângulo $\theta$ entre si} (são ângulos de lados mutuamente perpendiculares).
        \item Logo, o componente do peso colado ao eixo $y$ leva o cosseno:
        \[ P_y = P_\perp = P \cos\theta = mg \cos\theta \]
        \item O componente oposto (ao longo da rampa, puxando para baixo) leva o seno:
        \[ P_x = P_\parallel = P \sin\theta = mg \sin\theta \]
    \end{itemize}

    \item \textbf{O Teste Infalível dos Casos Limite (Para nunca trocar seno e cosseno):}
    Sempre faça o teste mental rápido antes de escrever a equação:
    \begin{itemize}
        \item Se o ângulo da rampa for $\theta \to 0^\circ$ (chão horizontal):
        $\sin 0^\circ = 0 \implies P_\parallel = 0$ (o corpo não desce sozinho);
        $\cos 0^\circ = 1 \implies P_\perp = mg$ (o peso inteiro aperta o chão).
        \item Se o ângulo da rampa for $\theta \to 90^\circ$ (parede vertical):
        $\sin 90^\circ = 1 \implies P_\parallel = mg$ (queda livre acelerada);
        $\cos 90^\circ = 0 \implies P_\perp = 0$ (o corpo não aperta a parede).
    \end{itemize}
    Se a sua fórmula der $mg\cos\theta$ para o componente que faz o corpo descer a rampa, você sabe imediatamente que errou, pois no chão plano ($\theta=0$) daria $mg$ para fazê-lo descer!

    \item \textbf{Força Externa Horizontal aplicada sobre Plano Inclinado:}
    Se uma força horizontal $\vec{F}_{\text{ext}}$ atua sobre um bloco na rampa, ela faz ângulo $\theta$ com a rampa. Decompondo nos eixos inclinados:
    \[ F_x = F_{\text{ext}} \cos\theta \quad \text{e} \quad F_y = -F_{\text{ext}} \sin\theta \]
    A componente perpendicular aumenta a compressão sobre a rampa: $n = mg\cos\theta + F_{\text{ext}}\sin\theta$.
\end{enumerate}

\vspace{0.4cm}

%%%%%%%%%%%%%%%%%%%%%%%%%%%%%%%%%%%%%%%%%%%%%%%%%%%%%%%%%%%%%%%%%%%%%%%%%%%%%%%
\section*{PARTE I: RESOLUÇÃO DAS QUESTÕES FÁCEIS (1 A 10)}
%%%%%%%%%%%%%%%%%%%%%%%%%%%%%%%%%%%%%%%%%%%%%%%%%%%%%%%%%%%%%%%%%%%%%%%%%%%%%%%

%------------------------------------------------------------------------------
\subsection*{Questão 1 [Sears 4.7]}
\textbf{Enunciado:} Se uma força resultante horizontal de $132\text{ N}$ é aplicada a uma pessoa com massa de $60\text{ kg}$ em repouso na beira de uma piscina, qual é a aceleração produzida?

\begin{dclbox}
\textbf{Corpo:} Pessoa ($m = 60\text{ kg}$).\\
\textbf{Forças verticais:} Peso $\vec{P}$ (Terra, para baixo, $P = mg$) e Normal $\vec{n}$ (chão, para cima).\\
Como não há movimento vertical: $\sum F_y = n - P = 0 \implies n = P = 588\text{ N}$.\\
\textbf{Forças horizontais:} Apenas a força resultante aplicada $F_{\text{res}} = 132\text{ N}$ no sentido $+x$.
\end{dclbox}

\noindent
\textbf{Resolução:}\\
Pela Segunda Lei de Newton na direção horizontal:
\[ \sum F_x = m a_x \implies F_{\text{res}} = m a \]
\[ a = \frac{F_{\text{res}}}{m} = \frac{132\text{ N}}{60\text{ kg}} = \mathbf{2,20\text{ m/s}^2} \]
\textbf{Resposta:} A aceleração produzida possui módulo igual a $\mathbf{2,20\text{ m/s}^2}$, na mesma direção e sentido da força aplicada.

%------------------------------------------------------------------------------
\vspace{0.3cm}
\subsection*{Questão 2 [Sears 4.8]}
\textbf{Enunciado:} Qual o módulo da força necessária para imprimir uma aceleração de $1,40\text{ m/s}^2$ em uma geladeira com massa de $135\text{ kg}$?

\begin{dclbox}
\textbf{Corpo:} Geladeira ($m = 135\text{ kg}$).\\
Na direção horizontal, para produzir a aceleração desejada $a = 1,40\text{ m/s}^2$, a força resultante horizontal deve satisfazer a 2ª Lei de Newton.
\end{dclbox}

\noindent
\textbf{Resolução:}\\
\[ F_{\text{res}} = m a = 135\text{ kg} \times 1,40\text{ m/s}^2 = \mathbf{189\text{ N}} \]
\textbf{Resposta:} O módulo da força necessária é $\mathbf{189\text{ N}}$.

%------------------------------------------------------------------------------
\vspace{0.3cm}
\subsection*{Questão 3 [Sears 4.9]}
\textbf{Enunciado:} Uma caixa está em repouso sobre um lago congelado, que é uma superfície horizontal sem atrito. Se um pescador aplica uma força horizontal de módulo $48,0\text{ N}$ sobre a caixa e produz uma aceleração de $3,00\text{ m/s}^2$, qual é a massa da caixa?

\begin{dclbox}
\textbf{Corpo:} Caixa.\\
Superfície sem atrito: $f_{\text{atrito}} = 0$.\\
Força horizontal aplicada: $F = 48,0\text{ N}$. Aceleração horizontal: $a = 3,00\text{ m/s}^2$.
\end{dclbox}

\noindent
\textbf{Resolução:}\\
Pela Segunda Lei de Newton:
\[ F = m a \implies m = \frac{F}{a} = \frac{48,0\text{ N}}{3,00\text{ m/s}^2} = \mathbf{16,0\text{ kg}} \]
\textbf{Resposta:} A massa da caixa é de $\mathbf{16,0\text{ kg}}$.

%------------------------------------------------------------------------------
\vspace{0.3cm}
\subsection*{Questão 4 [Sears 4.10]}
\textbf{Enunciado:} Um portuário aplica uma força horizontal constante de $80,0\text{ N}$ a um bloco de gelo sobre uma superfície horizontal lisa. A força de atrito é desprezível. O bloco parte do repouso e se move $11,0\text{ m}$ em $5,0\text{ s}$. a) Qual é a massa do bloco de gelo? b) Se o portuário parar de empurrar o bloco depois de $5,0\text{ s}$, qual será a distância percorrida pelo bloco nos $5,0\text{ s}$ posteriores?

\begin{dclbox}
\textbf{Corpo:} Bloco de gelo. Superfície horizontal sem atrito.\\
\textbf{Intervalo 1 ($0 \le t \le 5\text{ s}$):} Força constante $F = 80,0\text{ N}$, parte do repouso ($v_0 = 0$).\\
\textbf{Intervalo 2 ($5\text{ s} < t \le 10\text{ s}$):} Força cessa ($F = 0$).
\end{dclbox}

\noindent
\textbf{Resolução:}\\
a) Como a força é constante, o movimento no primeiro intervalo é uniformemente variado (MRUV):
\[ \Delta x = v_0 t + \frac{1}{2} a t^2 \implies 11,0 = 0 + \frac{1}{2} a (5,0)^2 \implies a = \frac{2 \times 11,0}{25,0} = 0,880\text{ m/s}^2 \]
Pela Segunda Lei de Newton:
\[ m = \frac{F}{a} = \frac{80,0\text{ N}}{0,880\text{ m/s}^2} = \mathbf{90,9\text{ kg}} \]
b) No instante $t = 5,0\text{ s}$, a velocidade atingida pelo bloco é:
\[ v_1 = v_0 + a t = 0 + (0,880\text{ m/s}^2)(5,0\text{ s}) = 4,40\text{ m/s} \]
Ao cessar a força ($F = 0$), pela Primeira Lei de Newton (Inércia), a força resultante torna-se nula e o bloco passa a se mover em Movimento Retilíneo Uniforme (MRU) com velocidade constante $v = 4,40\text{ m/s}$:
\[ \Delta x_2 = v_1 \times \Delta t_2 = (4,40\text{ m/s}) \times (5,0\text{ s}) = \mathbf{22,0\text{ m}} \]
\textbf{Resposta:} a) A massa é $\mathbf{90,9\text{ kg}}$. b) A distância percorrida nos $5,0\text{ s}$ seguintes é $\mathbf{22,0\text{ m}}$.

%------------------------------------------------------------------------------
\vspace{0.3cm}
\subsection*{Questão 5 [Sears 4.19]}
\textbf{Enunciado:} Na superfície de Io, uma das luas de Júpiter, a aceleração da gravidade é $g = 1,81\text{ m/s}^2$. Uma melancia pesa $44,0\text{ N}$ na superfície da Terra (onde $g = 9,80\text{ m/s}^2$). a) Qual sua massa na superfície da Terra? b) Qual sua massa e peso na superfície de Io?

\noindent
\textbf{Resolução:}\\
a) A massa $m$ é uma grandeza intrínseca do corpo:
\[ m = \frac{P_{\text{Terra}}}{g_{\text{Terra}}} = \frac{44,0\text{ N}}{9,80\text{ m/s}^2} = \mathbf{4,49\text{ kg}} \]
b) A massa não varia com o local, logo em Io ela continua sendo $\mathbf{4,49\text{ kg}}$.\\
O peso em Io depende da aceleração da gravidade local $g_{\text{Io}}$:
\[ P_{\text{Io}} = m \times g_{\text{Io}} = 4,49\text{ kg} \times 1,81\text{ m/s}^2 = \mathbf{8,13\text{ N}} \]
\textbf{Resposta:} a) $m = \mathbf{4,49\text{ kg}}$. b) Massa em Io: $\mathbf{4,49\text{ kg}}$; Peso em Io: $\mathbf{8,13\text{ N}}$.

%------------------------------------------------------------------------------
\vspace{0.3cm}
\subsection*{Questão 6 [Sears 4.20]}
\textbf{Enunciado:} A mochila de uma astronauta pesa $17,5\text{ N}$ quando ela está na superfície terrestre, mas somente $3,24\text{ N}$ na superfície de um asteroide. a) Qual é a aceleração da gravidade nesse asteroide? b) Qual é a massa da mochila no asteroide?

\noindent
\textbf{Resolução:}\\
Calculando a massa invariante da mochila a partir dos dados da Terra ($g_{\text{Terra}} = 9,80\text{ m/s}^2$):
\[ m = \frac{P_{\text{Terra}}}{g_{\text{Terra}}} = \frac{17,5\text{ N}}{9,80\text{ m/s}^2} = 1,786\text{ kg} \approx \mathbf{1,79\text{ kg}} \]
a) No asteroide, $P_{\text{ast}} = m \times g_{\text{ast}}$:
\[ g_{\text{ast}} = \frac{P_{\text{ast}}}{m} = \frac{3,24\text{ N}}{1,786\text{ kg}} = \mathbf{1,81\text{ m/s}^2} \]
b) A massa no asteroide é a mesma da Terra: $m = \mathbf{1,79\text{ kg}}$.\\
\textbf{Resposta:} a) $g_{\text{ast}} = \mathbf{1,81\text{ m/s}^2}$. b) $m = \mathbf{1,79\text{ kg}}$.

%------------------------------------------------------------------------------
\vspace{0.3cm}
\subsection*{Questão 7 [Sears 4.24]}
\textbf{Enunciado:} O piso de um elevador exerce uma força normal de $620\text{ N}$ de baixo para cima sobre um passageiro que pesa $650\text{ N}$. a) Quais são as reações dessas duas forças no sentido da Terceira Lei de Newton? b) O passageiro está sendo acelerado? Em caso afirmativo, determine o módulo, a direção e o sentido da aceleração.

\begin{dclbox}
\textbf{Corpo:} Passageiro.\\
Forças atuando sobre o passageiro:\\
1. Peso $\vec{P}$: atração gravitacional exercida pela Terra sobre o passageiro ($650\text{ N}$, vertical para baixo).\\
2. Força normal $\vec{n}$: força de contato exercida pelo piso sobre o passageiro ($620\text{ N}$, vertical para cima).
\end{dclbox}

\noindent
\textbf{Resolução:}\\
a) Pela Terceira Lei de Newton (Ação e Reação):
\begin{itemize}
    \item A reação à força normal $\vec{n}$ é a \textbf{força normal que os pés do passageiro exercem sobre o piso do elevador} (módulo $620\text{ N}$, vertical para baixo).
    \item A reação ao peso $\vec{P}$ é a \textbf{força de atração gravitacional que o passageiro exerce sobre o planeta Terra} (módulo $650\text{ N}$, vertical para cima, aplicada no centro da Terra).
\end{itemize}
b) Força resultante sobre o passageiro no eixo vertical ($+y$ para cima):
\[ \sum F_y = n - P = 620\text{ N} - 650\text{ N} = -30\text{ N} \]
Como a força resultante não é nula, \textbf{o passageiro está acelerado para baixo}.\\
Massa do passageiro: $m = \frac{P}{g} = \frac{650\text{ N}}{9,80\text{ m/s}^2} = 66,33\text{ kg}$.\\
Aceleração:
\[ a_y = \frac{\sum F_y}{m} = \frac{-30\text{ N}}{66,33\text{ kg}} = \mathbf{-0,452\text{ m/s}^2} \]
\textbf{Resposta:} Sim, acelerado com módulo $\mathbf{0,452\text{ m/s}^2}$, na direção vertical com sentido \textbf{para baixo}.

%------------------------------------------------------------------------------
\vspace{0.3cm}
\subsection*{Questão 8 [Sears 5.1]}
\textbf{Enunciado:} Dois pesos de $25,0\text{ N}$ estão suspensos nas extremidades opostas de uma corda leve que passa sobre uma polia sem atrito. O centro da polia está ligado a uma corrente presa ao teto. a) Qual é a tensão na corda? b) Qual é a tensão na corrente presa ao teto?

\begin{dclbox}
\textbf{DCL de cada bloco de peso $P = 25,0\text{ N}$:} Tração $T$ da corda para cima e Peso $P$ para baixo.\\
\textbf{DCL da polia (sem massa):} Tração $T_{\text{corrente}}$ da corrente para cima, e dois segmentos da corda puxando para baixo, cada um com tração $T$.
\end{dclbox}

\noindent
\textbf{Resolução:}\\
a) Como os dois pesos são idênticos ($25,0\text{ N}$), o sistema está em equilíbrio estático ($a = 0$):
\[ \sum F_y = T - P = 0 \implies T = P = \mathbf{25,0\text{ N}} \]
b) Para a polia em equilíbrio (massa desprezível):
\[ \sum F_y = T_{\text{corrente}} - 2T = 0 \implies T_{\text{corrente}} = 2 \times 25,0\text{ N} = \mathbf{50,0\text{ N}} \]
\textbf{Resposta:} a) Tensão na corda: $\mathbf{25,0\text{ N}}$. b) Tensão na corrente: $\mathbf{50,0\text{ N}}$.

%------------------------------------------------------------------------------
\vspace{0.3cm}
\subsection*{Questão 9 [Sears 5.11]}
\textbf{Enunciado:} Um homem empurra um piano de $180\text{ kg}$, de modo que ele desliza com velocidade constante para baixo de uma rampa inclinada de $11,0^\circ$ acima da horizontal. Despreze o atrito que atua sobre o piano. Calcule o módulo da força aplicada pelo homem se ela for: a) paralela ao plano inclinado; b) paralela ao piso horizontal.

\begin{dclbox}
\textbf{Corpo:} Piano ($m = 180\text{ kg}$, $P = mg = 180 \times 9,80 = 1764\text{ N}$).\\
\textbf{Velocidade constante:} $\vec{a} = 0$ (equilíbrio dinâmico).\\
\textbf{Eixos:} $x$ paralelo à rampa (apontando rampa acima), $y$ perpendicular à rampa.\\
Componentes do peso: $P_\parallel = P\sin(11,0^\circ)$ para baixo da rampa; $P_\perp = P\cos(11,0^\circ)$ para dentro da rampa.
\end{dclbox}

\noindent
\textbf{Resolução:}\\
a) Se a força $\vec{F}$ é aplicada paralela à rampa (empurrando rampa acima para segurar a descida com velocidade constante):
\[ \sum F_x = F - P\sin(11,0^\circ) = 0 \implies F = mg\sin(11,0^\circ) \]
\[ F = 1764\text{ N} \times \sin(11,0^\circ) = 1764 \times 0,1908 = \mathbf{337\text{ N}} \]
b) Se a força $\vec{F}$ é aplicada horizontalmente (empurrando contra a rampa):\\
Como a rampa faz $11,0^\circ$ com a horizontal, o vetor horizontal $\vec{F}$ faz ângulo de $11,0^\circ$ com a linha da rampa.
\begin{itemize}
    \item Componente ao longo da rampa (para cima): $F \cos(11,0^\circ)$.
    \item Componente perpendicular à rampa (apertando a rampa): $F \sin(11,0^\circ)$.
\end{itemize}
Equilíbrio ao longo da rampa ($\sum F_x = 0$):
\[ F\cos(11,0^\circ) - P\sin(11,0^\circ) = 0 \implies F = P \frac{\sin(11,0^\circ)}{\cos(11,0^\circ)} = mg \tan(11,0^\circ) \]
\[ F = 1764\text{ N} \times \tan(11,0^\circ) = 1764 \times 0,19438 = \mathbf{343\text{ N}} \]
\textbf{Resposta:} a) $\mathbf{337\text{ N}}$. b) $\mathbf{343\text{ N}}$.

%------------------------------------------------------------------------------
\vspace{0.3cm}
\subsection*{Questão 10 [Sears 5.28]}
\textbf{Enunciado:} Uma caixa de $25,0\text{ kg}$ com livros repousa sobre uma mesa horizontal. O coeficiente de atrito estático entre a caixa e a mesa é $0,40$ e o coeficiente de atrito cinético é $0,20$. a) Qual é o peso da caixa e o módulo da força normal exercida pela mesa sobre ela? b) Qual força horizontal é necessária para iniciar o movimento da caixa? c) Qual força horizontal é necessária para manter a caixa em movimento com velocidade constante uma vez iniciado o deslizamento?

\begin{dclbox}
\textbf{Corpo:} Caixa ($m = 25,0\text{ kg}$). Superfície plana horizontal.\\
Vertical: Peso $P = mg$ (para baixo) e Normal $n$ (para cima).\\
Horizontal: Força aplicada $F$ (para a direita) e Atrito $f_{\text{atrito}}$ (para a esquerda).
\end{dclbox}

\noindent
\textbf{Resolução:}\\
a) $P = mg = 25,0\text{ kg} \times 9,80\text{ m/s}^2 = \mathbf{245\text{ N}}$.\\
Como não há aceleração vertical: $\sum F_y = n - P = 0 \implies n = P = \mathbf{245\text{ N}}$.\\
b) Para iniciar o movimento, a força horizontal deve superar a força de atrito estático máxima:
\[ F_{\text{início}} = f_{e,\text{max}} = \mu_e \times n = 0,40 \times 245\text{ N} = \mathbf{98,0\text{ N}} \]
c) Em movimento de deslizamento uniforme ($a = 0$), a força necessária equilibra o atrito cinético:
\[ F_{\text{mru}} = f_c = \mu_c \times n = 0,20 \times 245\text{ N} = \mathbf{49,0\text{ N}} \]
\textbf{Resposta:} a) $P = 245\text{ N}$ e $n = 245\text{ N}$. b) $F \ge \mathbf{98,0\text{ N}}$. c) $F = \mathbf{49,0\text{ N}}$.

\newpage
%%%%%%%%%%%%%%%%%%%%%%%%%%%%%%%%%%%%%%%%%%%%%%%%%%%%%%%%%%%%%%%%%%%%%%%%%%%%%%%
\section*{PARTE II: RESOLUÇÃO DAS QUESTÕES MÉDIAS (11 A 25)}
%%%%%%%%%%%%%%%%%%%%%%%%%%%%%%%%%%%%%%%%%%%%%%%%%%%%%%%%%%%%%%%%%%%%%%%%%%%%%%%

%------------------------------------------------------------------------------
\subsection*{Questão 11 [Sears 4.5]}
\textbf{Enunciado:} Dois cachorros puxam horizontalmente cordas amarradas a um poste; o ângulo entre as duas cordas é igual a $60,0^\circ$. Se o cachorro A exerce uma força de $270\text{ N}$ e o cachorro B exerce uma força de $300\text{ N}$, ache o módulo da força resultante sobre o poste e o ângulo que ela forma com a corda puxada pelo cachorro A.

\begin{dclbox}
\textbf{Corpo:} Ponto de amarração no poste.\\
Escolha dos eixos: orientamos o eixo $+x$ ao longo da corda A.\\
$\vec{F}_A = (270\text{ N})\hat{\imath} + (0)\hat{\jmath}$.\\
$\vec{F}_B$ faz $60,0^\circ$ com $+x$: $F_{Bx} = 300\cos(60^\circ)$, $F_{By} = 300\sin(60^\circ)$.
\end{dclbox}

\noindent
\textbf{Resolução:}\\
Calculando os componentes cartesianos da força resultante $\vec{R} = \vec{F}_A + \vec{F}_B$:
\[ R_x = F_{Ax} + F_{Bx} = 270 + 300 \cos(60,0^\circ) = 270 + 300(0,500) = 270 + 150 = 420\text{ N} \]
\[ R_y = F_{Ay} + F_{By} = 0 + 300 \sin(60,0^\circ) = 300(0,8660) = 259,8\text{ N} \]
Módulo da força resultante:
\[ R = \sqrt{R_x^2 + R_y^2} = \sqrt{(420)^2 + (259,8)^2} = \sqrt{176\,400 + 67\,496} = \sqrt{243\,896} = \mathbf{494\text{ N}} \]
Direção (ângulo $\theta$ com a corda A):
\[ \tan\theta = \frac{R_y}{R_x} = \frac{259,8}{420} = 0,6186 \implies \theta = \arctan(0,6186) = \mathbf{31,7^\circ} \]
\textbf{Resposta:} A força resultante possui módulo de $\mathbf{494\text{ N}}$ e forma um ângulo de $\mathbf{31,7^\circ}$ com a corda do cachorro A.

%------------------------------------------------------------------------------
\vspace{0.3cm}
\subsection*{Questão 12 [Sears 4.11]}
\textbf{Enunciado:} Um disco de hóquei com massa de $0,160\text{ kg}$ está em repouso na origem ($x = 0$) em uma superfície horizontal sem atrito da pista. No instante $t = 0$, um jogador aplica uma força constante de $0,250\text{ N}$ sobre o disco, paralela ao eixo $Ox$. Ele continua a aplicar essa força até o instante $t = 2,00\text{ s}$. a) Qual é a posição e a velocidade do disco no instante $t = 2,00\text{ s}$? b) Se a mesma força for aplicada novamente no instante $t = 5,00\text{ s}$ no mesmo sentido, qual será a posição e a velocidade do disco no instante $t = 7,00\text{ s}$?

\noindent
\textbf{Resolução:}\\
Aceleração quando a força de $0,250\text{ N}$ está aplicada:
\[ a = \frac{F}{m} = \frac{0,250\text{ N}}{0,160\text{ kg}} = 1,5625\text{ m/s}^2 \]
a) No intervalo de $t = 0$ a $t = 2,00\text{ s}$ ($v_0 = 0$, $x_0 = 0$):
\[ v(2) = a t = (1,5625\text{ m/s}^2)(2,00\text{ s}) = \mathbf{3,13\text{ m/s}} \]
\[ x(2) = \frac{1}{2} a t^2 = \frac{1}{2} (1,5625)(2,00)^2 = \mathbf{3,13\text{ m}} \]
b) No intervalo de $t = 2,00\text{ s}$ a $t = 5,00\text{ s}$, a força é zero ($a = 0$, MRU por $\Delta t = 3,00\text{ s}$):
\[ v(5) = 3,125\text{ m/s} \]
\[ x(5) = x(2) + v(2)\Delta t = 3,125 + (3,125\text{ m/s})(3,00\text{ s}) = 3,125 + 9,375 = 12,50\text{ m} \]
No intervalo de $t = 5,00\text{ s}$ a $t = 7,00\text{ s}$, a força é reaplicada ($a = 1,5625\text{ m/s}^2$ durante $\Delta t = 2,00\text{ s}$):
\[ v(7) = v(5) + a \Delta t = 3,125 + (1,5625)(2,00) = 3,125 + 3,125 = \mathbf{6,25\text{ m/s}} \]
\[ x(7) = x(5) + v(5)\Delta t + \frac{1}{2} a (\Delta t)^2 = 12,50 + (3,125)(2,00) + \frac{1}{2}(1,5625)(2,00)^2 = 12,50 + 6,25 + 3,125 = \mathbf{21,9\text{ m}} \]
\textbf{Resposta:} a) $x = \mathbf{3,13\text{ m}}$ e $v = \mathbf{3,13\text{ m/s}}$. b) $x = \mathbf{21,9\text{ m}}$ e $v = \mathbf{6,25\text{ m/s}}$.

%------------------------------------------------------------------------------
\vspace{0.3cm}
\subsection*{Questão 13 [Sears 4.17]}
\textbf{Enunciado:} Uma caixa com massa de $5,0\text{ kg}$ é puxada por uma corda ao longo de uma superfície horizontal lisa sem atrito. A tensão na corda é de $40,0\text{ N}$ e ela faz um ângulo de $30,0^\circ$ acima da horizontal. a) Qual é a aceleração horizontal da caixa? b) Qual é o módulo da força normal que a superfície exerce sobre a caixa?

\begin{dclbox}
\textbf{Corpo:} Caixa ($m = 5,0\text{ kg}$).\\
Forças que atuam sobre a caixa:\\
1. Peso $\vec{P}$: vertical para baixo, $P = mg = 5,0 \times 9,80 = 49,0\text{ N}$.\\
2. Tração $\vec{T}$: $T = 40,0\text{ N}$ a $30,0^\circ$ acima da horizontal.\\
   Componentes: $T_x = T\cos(30^\circ) = 40,0(0,8660) = 34,64\text{ N}$; $T_y = T\sin(30^\circ) = 40,0(0,500) = 20,0\text{ N}$.\\
3. Força normal $\vec{n}$: vertical para cima (chão).
\end{dclbox}

\noindent
\textbf{Resolução:}\\
a) Na direção horizontal (sem atrito):
\[ \sum F_x = T_x = m a_x \implies a_x = \frac{T_x}{m} = \frac{34,64\text{ N}}{5,0\text{ kg}} = \mathbf{6,93\text{ m/s}^2} \]
b) Na direção vertical (sem movimento vertical, $\sum F_y = 0$):
\[ \sum F_y = n + T_y - P = 0 \implies n = mg - T_y = 49,0\text{ N} - 20,0\text{ N} = \mathbf{29,0\text{ N}} \]
\textbf{Resposta:} a) $a = \mathbf{6,93\text{ m/s}^2}$. b) $n = \mathbf{29,0\text{ N}}$.

%------------------------------------------------------------------------------
\vspace{0.3cm}
\subsection*{Questão 14 [Sears 4.25]}
\textbf{Enunciado:} Uma estudante com massa de $45\text{ kg}$ pula de um trampolim elevado. Considerando a massa da Terra como $6,0 \times 10^{24}\text{ kg}$, qual é a aceleração da Terra no sentido da estudante quando ela se acelera no sentido da Terra com $9,80\text{ m/s}^2$? Suponha que a força resultante sobre a Terra seja a força gravitacional que a estudante exerce sobre ela.

\noindent
\textbf{Resolução:}\\
Pela Terceira Lei de Newton, a força gravitacional que a Terra exerce sobre a estudante (seu peso) é igual em módulo e oposta em sentido à força gravitacional que a estudante exerce sobre a Terra:
\[ F_{\text{estudante na Terra}} = F_{\text{Terra na estudante}} = m_e \times g = 45\text{ kg} \times 9,80\text{ m/s}^2 = 441\text{ N} \]
Pela Segunda Lei de Newton aplicada ao planeta Terra:
\[ a_{\text{Terra}} = \frac{F}{M_{\text{Terra}}} = \frac{441\text{ N}}{6,0 \times 10^{24}\text{ kg}} = \mathbf{7,35 \times 10^{-23}\text{ m/s}^2} \]
\textbf{Resposta:} A aceleração da Terra em direção à estudante é de aproximadamente $\mathbf{7,4 \times 10^{-23}\text{ m/s}^2}$.

%------------------------------------------------------------------------------
\vspace{0.3cm}
\subsection*{Questão 15 [Sears 4.36]}
\textbf{Enunciado:} Um estudante de física com $550\text{ N}$ de peso sobe em uma balança de mola apoiada no chão de um elevador. Conforme o elevador se move, a balança registra uma leitura de peso aparente de $450\text{ N}$. a) Determine o módulo, a direção e o sentido da aceleração do elevador. b) Se a balança registrar $670\text{ N}$, qual será a aceleração do elevador? c) Se a leitura da balança for nula, o que se pode concluir sobre o movimento do elevador?

\begin{dclbox}
\textbf{Corpo:} Estudante.\\
A balança mede a \textbf{força normal $n$} exercida sobre os pés da pessoa.\\
Massa da pessoa: $m = \frac{P}{g} = \frac{550\text{ N}}{9,80\text{ m/s}^2} = 56,12\text{ kg}$.\\
Equação do movimento no eixo vertical ($+y$ para cima): $\sum F_y = n - P = m a$.
\end{dclbox}

\noindent
\textbf{Resolução:}\\
a) Para leitura $n = 450\text{ N}$:
\[ a = \frac{n - P}{m} = \frac{450 - 550}{56,12} = \frac{-100\text{ N}}{56,12\text{ kg}} = \mathbf{-1,78\text{ m/s}^2} \]
Módulo $\mathbf{1,78\text{ m/s}^2}$, vertical para baixo.\\
b) Para leitura $n = 670\text{ N}$:
\[ a = \frac{670 - 550}{56,12} = \frac{+120\text{ N}}{56,12\text{ kg}} = \mathbf{+2,14\text{ m/s}^2} \]
Módulo $\mathbf{2,14\text{ m/s}^2}$, vertical para cima.\\
c) Se $n = 0$: $a = \frac{0 - 550}{56,12} = -9,80\text{ m/s}^2 = -g$. O elevador está em \textbf{queda livre} (os cabos romperam ou o elevador acelera para baixo com $g$).\\
\textbf{Resposta:} a) $1,78\text{ m/s}^2$ para baixo. b) $2,14\text{ m/s}^2$ para cima. c) Queda livre ($a = g$ para baixo).

%------------------------------------------------------------------------------
\vspace{0.3cm}
\subsection*{Questão 16 [Sears 4.39]}
\textbf{Enunciado:} Dois blocos de massas $m_1 = 4,0\text{ kg}$ e $m_2 = 6,0\text{ kg}$ estão apoiados lado a lado em contato mútuo sobre uma superfície horizontal sem atrito. Uma força horizontal constante $F = 30,0\text{ N}$ é aplicada ao bloco $m_1$, empurrando-o contra o bloco $m_2$. a) Desenhe o diagrama do corpo livre para cada um dos blocos. b) Qual é a aceleração do sistema formado pelos dois blocos? c) Qual é o módulo da força de contato exercida entre os dois blocos?

\begin{dclbox}
\textbf{DCL do bloco 1 ($m_1$):} Peso $P_1$, Normal $n_1$, Força aplicada $F$ (para a direita), Força de contato $F_{c}$ (exercida por $m_2$ sobre $m_1$, para a esquerda).\\
\textbf{DCL do bloco 2 ($m_2$):} Peso $P_2$, Normal $n_2$, Força de contato $F_{c}$ (exercida por $m_1$ sobre $m_2$, para a direita).
\end{dclbox}

\noindent
\textbf{Resolução:}\\
b) Considerando o sistema conjunto $(m_1 + m_2)$, a força de contato é interna. A força resultante externa é apenas $F$:
\[ a = \frac{F}{m_1 + m_2} = \frac{30,0\text{ N}}{4,0\text{ kg} + 6,0\text{ kg}} = \frac{30,0\text{ N}}{10,0\text{ kg}} = \mathbf{3,00\text{ m/s}^2} \]
c) Aplicando a Segunda Lei isoladamente sobre o bloco 2:
\[ \sum F_{x,2} = F_c = m_2 a = 6,0\text{ kg} \times 3,00\text{ m/s}^2 = \mathbf{18,0\text{ N}} \]
(Verificação no bloco 1: $F - F_c = m_1 a \implies 30,0 - 18,0 = 12,0\text{ N} = 4,0 \times 3,00$, consistente!).\\
\textbf{Resposta:} b) $a = \mathbf{3,00\text{ m/s}^2}$. c) Força de contato: $\mathbf{18,0\text{ N}}$.

%------------------------------------------------------------------------------
\vspace{0.3cm}
\subsection*{Questão 17 [Sears 5.3]}
\textbf{Enunciado:} Um quadro está suspenso em uma parede por meio de dois fios presos aos seus cantos superiores. Cada fio forma um ângulo de $30,0^\circ$ com a vertical. Se a massa total do quadro for de $1,50\text{ kg}$, calcule a tensão em cada um dos fios.

\begin{dclbox}
\textbf{Corpo:} Quadro ($m = 1,50\text{ kg}$, $P = mg = 1,50 \times 9,80 = 14,7\text{ N}$).\\
Dois fios simétricos com tensão $T$, cada um formando ângulo $\theta = 30,0^\circ$ com a vertical.\\
Componente vertical de cada fio: $T_y = T \cos(30,0^\circ)$.\\
Componentes horizontais: $T \sin(30,0^\circ)$ opostos que se cancelam.
\end{dclbox}

\noindent
\textbf{Resolução:}\\
Condição de equilíbrio estático no eixo vertical ($\sum F_y = 0$):
\[ 2 T \cos(30,0^\circ) - P = 0 \implies 2 T \cos(30,0^\circ) = mg \]
\[ T = \frac{mg}{2 \cos(30,0^\circ)} = \frac{14,7\text{ N}}{2 \times 0,8660} = \frac{14,7}{1,732} = \mathbf{8,49\text{ N}} \]
\textbf{Resposta:} A tensão em cada fio é de $\mathbf{8,49\text{ N}}$.

%------------------------------------------------------------------------------
\vspace{0.3cm}
\subsection*{Questão 18 [Sears 5.9]}
\textbf{Enunciado:} Um carro de $1000\text{ kg}$ está sobre uma rampa sem atrito inclinada de $20,0^\circ$ em relação à horizontal, sendo sustentado por um cabo de aço paralelo à rampa. a) Se o carro sobe a rampa com velocidade constante de $15,0\text{ m/s}$, qual é a tensão no cabo? b) Se o carro sobe a rampa acelerando a uma taxa de $2,00\text{ m/s}^2$, qual é a tensão no cabo?

\begin{dclbox}
\textbf{Corpo:} Carro ($m = 1000\text{ kg}$, $P = mg = 9800\text{ N}$).\\
\textbf{Eixos inclinados:} $+x$ rampa acima; $+y$ perpendicular à rampa.\\
Componente do peso puxando rampa abaixo: $P_\parallel = mg\sin(20,0^\circ) = 9800 \times 0,3420 = 3352\text{ N}$.
\end{dclbox}

\noindent
\textbf{Resolução:}\\
a) Subida com velocidade constante ($a = 0$):
\[ \sum F_x = T - mg\sin(20,0^\circ) = 0 \implies T = 3352\text{ N} \approx \mathbf{3,35 \times 10^3\text{ N}} \]
b) Subida acelerada com $a = 2,00\text{ m/s}^2$:
\[ \sum F_x = T - mg\sin(20,0^\circ) = m a \]
\[ T = mg\sin(20,0^\circ) + m a = 3352\text{ N} + (1000\text{ kg})(2,00\text{ m/s}^2) = 3352 + 2000 = \mathbf{5,35 \times 10^3\text{ N}} \]
\textbf{Resposta:} a) $T = \mathbf{3,35 \times 10^3\text{ N}}$. b) $T = \mathbf{5,35 \times 10^3\text{ N}}$.

%------------------------------------------------------------------------------
\vspace{0.3cm}
\subsection*{Questão 19 [Sears 5.17]}
\textbf{Enunciado:} Um bloco de gelo de $8,00\text{ kg}$, liberado a partir do repouso no topo de uma rampa sem atrito com $1,50\text{ m}$ de comprimento, desliza para baixo, alcançando uma velocidade escalar de $2,50\text{ m/s}$ na base da rampa. a) Qual é o ângulo de inclinação da rampa com a horizontal? b) Qual seria a velocidade escalar do bloco na base se ele fosse lançado para cima a partir da base ao longo da mesma rampa com velocidade inicial de $2,50\text{ m/s}$?

\noindent
\textbf{Resolução:}\\
a) Pela Equação de Torricelli ($v_0 = 0$, $s = 1,50\text{ m}$, $v = 2,50\text{ m/s}$):
\[ v^2 = v_0^2 + 2 a s \implies a = \frac{v^2}{2 s} = \frac{(2,50)^2}{2(1,50)} = \frac{6,25}{3,00} = 2,083\text{ m/s}^2 \]
No plano inclinado sem atrito, a aceleração ao longo da rampa é produzida exclusivamente pela componente do peso:
\[ a = g \sin\theta \implies \sin\theta = \frac{a}{g} = \frac{2,083\text{ m/s}^2}{9,80\text{ m/s}^2} = 0,2126 \]
\[ \theta = \arcsin(0,2126) = \mathbf{12,3^\circ} \]
b) Como não há atrito nem outras forças não conservativas, a energia mecânica se conserva. O bloco desacelera até o repouso no topo e retorna descendo com aceleração idêntica. Ao atingir a base novamente, sua velocidade escalar será exatamente o mesmo valor inicial: $\mathbf{2,50\text{ m/s}}$.\\
\textbf{Resposta:} a) $\theta = \mathbf{12,3^\circ}$. b) $v = \mathbf{2,50\text{ m/s}}$.

%------------------------------------------------------------------------------
\vspace{0.3cm}
\subsection*{Questão 20 [Sears 5.19]}
\textbf{Enunciado:} \textbf{Máquina de Atwood.} Uma carga de tijolos com massa de $15,0\text{ kg}$ é suspensa pela extremidade de uma corda leve que passa sobre uma pequena polia fixa sem atrito. Um contrapeso de $28,0\text{ kg}$ está preso na outra extremidade da corda. O sistema é liberado a partir do repouso. a) Desenhe o diagrama do corpo livre para a carga de tijolos e outro para o contrapeso. b) Qual é o módulo da aceleração do sistema? c) Qual é a tensão na corda durante o movimento acelerado?

\begin{dclbox}
\textbf{Tijolos ($m_1 = 15,0\text{ kg}$):} aceleram para cima ($+a$).\\
Forças: Tração $T$ (para cima), Peso $P_1 = m_1 g$ (para baixo). Equação: $T - m_1 g = m_1 a$.\\
\textbf{Contrapeso ($m_2 = 28,0\text{ kg}$):} acelera para baixo ($-a$).\\
Forças: Tração $T$ (para cima), Peso $P_2 = m_2 g$ (para baixo). Equação: $m_2 g - T = m_2 a$.
\end{dclbox}

\noindent
\textbf{Resolução:}\\
b) Somando as equações de movimento dos dois corpos:
\[ (T - m_1 g) + (m_2 g - T) = m_1 a + m_2 a \implies (m_2 - m_1)g = (m_1 + m_2)a \]
\[ a = \frac{m_2 - m_1}{m_1 + m_2} g = \frac{28,0 - 15,0}{28,0 + 15,0} \times 9,80 = \frac{13,0}{43,0} \times 9,80 = \mathbf{2,96\text{ m/s}^2} \]
c) Substituindo na equação dos tijolos:
\[ T = m_1 (g + a) = 15,0\text{ kg} \times (9,80 + 2,96)\text{ m/s}^2 = 15,0 \times 12,76 = \mathbf{191\text{ N}} \]
\textbf{Resposta:} b) $a = \mathbf{2,96\text{ m/s}^2}$. c) Tensão na corda: $\mathbf{191\text{ N}}$ (maior que o peso dos tijolos de $147\text{ N}$ e menor que o peso do contrapeso de $274\text{ N}$).

%------------------------------------------------------------------------------
\vspace{0.3cm}
\subsection*{Questão 21 [Sears 5.23]}
\textbf{Enunciado:} Um tobogã com massa de $75,0\text{ kg}$ desce uma encosta coberta de neve inclinada de $25,0^\circ$ com a horizontal. O coeficiente de atrito cinético entre os patins do tobogã e a neve é $\mu_c = 0,12$. a) Desenhe o diagrama do corpo livre para o tobogã. b) Calcule o módulo da força normal exercida pela encosta sobre o tobogã. c) Calcule a aceleração do tobogã ao longo da encosta.

\begin{dclbox}
\textbf{Eixos inclinados:} $+x$ descendo a rampa; $+y$ perpendicular para fora da rampa.\\
1. Força Peso $\vec{P}$ ($P = mg = 75,0 \times 9,80 = 735\text{ N}$):\\
   $P_x = mg\sin(25,0^\circ) = 735 \times 0,4226 = 310,6\text{ N}$;\\
   $P_y = -mg\cos(25,0^\circ) = -735 \times 0,9063 = -666,1\text{ N}$.\\
2. Força Normal $\vec{n}$: no sentido $+y$.\\
3. Força de atrito cinético $\vec{f}_c$: no sentido $-x$ (oposta ao deslizamento).
\end{dclbox}

\noindent
\textbf{Resolução:}\\
b) No eixo perpendicular ($y$), não há movimento:
\[ \sum F_y = n - mg\cos(25,0^\circ) = 0 \implies n = 735 \times \cos(25,0^\circ) = \mathbf{666\text{ N}} \]
c) Força de atrito cinético:
\[ f_c = \mu_c n = 0,12 \times 666,1\text{ N} = 79,9\text{ N} \]
Segunda Lei de Newton no eixo paralelo ($x$):
\[ \sum F_x = mg\sin(25,0^\circ) - f_c = m a \]
\[ a = g\sin(25,0^\circ) - \mu_c g\cos(25,0^\circ) = 9,80(0,4226 - 0,12 \times 0,9063) = 9,80(0,4226 - 0,1088) = \mathbf{3,08\text{ m/s}^2} \]
\textbf{Resposta:} b) $n = \mathbf{666\text{ N}}$. c) $a = \mathbf{3,08\text{ m/s}^2}$.

%------------------------------------------------------------------------------
\vspace{0.3cm}
\subsection*{Questão 22 [Sears 5.31]}
\textbf{Enunciado:} Uma caixa de ferramentas de $45,0\text{ kg}$ está em repouso sobre um piso horizontal. Você exerce sobre ela uma força de puxar horizontal que aumenta gradualmente e observa que a caixa só começa a se mover quando a sua força ultrapassa $313\text{ N}$. A partir daí, você deve reduzir sua força de puxar para $208\text{ N}$ para mantê-la em movimento com velocidade constante de $25,0\text{ cm/s}$. a) Quais são os coeficientes de atrito estático ($\mu_e$) e cinético ($\mu_c$) entre a caixa e o piso? b) Que força horizontal de puxar você deve exercer para imprimir à caixa uma aceleração de $1,10\text{ m/s}^2$? c) Suponha que você realize essa experiência na Lua ($g = 1,62\text{ m/s}^2$). Qual seria a força necessária para iniciar o movimento?

\noindent
\textbf{Resolução:}\\
Força normal no piso horizontal: $n = mg = 45,0\text{ kg} \times 9,80\text{ m/s}^2 = 441\text{ N}$.\\
a) No limiar de escorregamento: $f_{e,\text{max}} = 313\text{ N}$:
\[ \mu_e = \frac{f_{e,\text{max}}}{n} = \frac{313\text{ N}}{441\text{ N}} = \mathbf{0,710} \]
Em velocidade constante ($a = 0$): $f_c = 208\text{ N}$:
\[ \mu_c = \frac{f_c}{n} = \frac{208\text{ N}}{441\text{ N}} = \mathbf{0,472} \]
b) Para acelerar com $a = 1,10\text{ m/s}^2$:
\[ F - f_c = m a \implies F = f_c + m a = 208\text{ N} + (45,0\text{ kg})(1,10\text{ m/s}^2) = 208 + 49,5 = \mathbf{258\text{ N}} \]
c) Na Lua, a força normal diminui para: $n_{\text{lua}} = m g_{\text{lua}} = 45,0 \times 1,62 = 72,9\text{ N}$.\\
Como os coeficientes de atrito dependem apenas dos materiais em contato:
\[ F_{\text{início,lua}} = \mu_e n_{\text{lua}} = 0,710 \times 72,9\text{ N} = \mathbf{51,8\text{ N}} \]
\textbf{Resposta:} a) $\mu_e = \mathbf{0,710}$ e $\mu_c = \mathbf{0,472}$. b) $F = \mathbf{258\text{ N}}$. c) $F_{\text{lua}} = \mathbf{51,8\text{ N}}$.

%------------------------------------------------------------------------------
\vspace{0.3cm}
\subsection*{Questão 23 [Sears 5.33]}
\textbf{Enunciado:} Um engradado de $48,0\text{ kg}$ está sobre um piso horizontal rugoso com coeficiente de atrito cinético $\mu_c = 0,30$. Você puxa o engradado por meio de uma corda que faz um ângulo de $30,0^\circ$ acima da horizontal. Se a tensão na corda for de $250\text{ N}$: a) Calcule o módulo da força normal exercida pelo piso sobre o engradado. b) Calcule a força de atrito cinético exercida sobre o engradado. c) Calcule a aceleração horizontal do engradado.

\begin{dclbox}
\textbf{Corpo:} Engradado ($m = 48,0\text{ kg}$, $P = mg = 48,0 \times 9,80 = 470,4\text{ N}$).\\
Decomposição da tração $\vec{T}$ ($250\text{ N}$ a $30,0^\circ$):\\
$T_x = 250 \cos(30^\circ) = 216,5\text{ N}$; $T_y = 250 \sin(30^\circ) = 125,0\text{ N}$.
\end{dclbox}

\noindent
\textbf{Resolução:}\\
a) Equilíbrio vertical ($\sum F_y = 0$):
\[ n + T_y - P = 0 \implies n = P - T_y = 470,4\text{ N} - 125,0\text{ N} = \mathbf{345\text{ N}} \]
b) Força de atrito cinético:
\[ f_c = \mu_c n = 0,30 \times 345,4\text{ N} = \mathbf{104\text{ N}} \]
c) Segunda Lei de Newton na horizontal:
\[ \sum F_x = T_x - f_c = m a \implies a = \frac{216,5\text{ N} - 103,6\text{ N}}{48,0\text{ kg}} = \frac{112,9}{48,0} = \mathbf{2,35\text{ m/s}^2} \]
\textbf{Resposta:} a) $n = \mathbf{345\text{ N}}$. b) $f_c = \mathbf{104\text{ N}}$. c) $a = \mathbf{2,35\text{ m/s}^2}$.

%------------------------------------------------------------------------------
\vspace{0.3cm}
\subsection*{Questão 24 [Sears 5.43]}
\textbf{Enunciado:} Uma curva plana em uma rodovia possui raio de curvatura igual a $220\text{ m}$. Em um dia chuvoso, o coeficiente de atrito estático entre os pneus de um veículo e o asfalto é $\mu_e = 0,40$. Qual é a velocidade máxima com que um automóvel pode realizar essa curva sem derrapar lateralmente (expresse a resposta em m/s e em km/h)?

\begin{dclbox}
\textbf{Curva plana (sem inclinação lateral):}\\
A única força horizontal que aponta para o centro da curva e fornece a aceleração centrípeta é a \textbf{força de atrito estático lateral} $\vec{f}_e$.\\
Vertical: $n = mg$.
\end{dclbox}

\noindent
\textbf{Resolução:}\\
Para não derrapar:
\[ f_e \le f_{e,\text{max}} = \mu_e n = \mu_e mg \]
A força centrípeta necessária para manter a trajetória de raio $R$ com velocidade $v$ é:
\[ F_{\text{rad}} = m \frac{v^2}{R} \]
No limite máximo de velocidade:
\[ m \frac{v_{\text{max}}^2}{R} = \mu_e mg \implies v_{\text{max}} = \sqrt{\mu_e g R} \]
\[ v_{\text{max}} = \sqrt{0,40 \times 9,80\text{ m/s}^2 \times 220\text{ m}} = \sqrt{862,4} = \mathbf{29,4\text{ m/s}} \]
Convertendo para km/h:
\[ v_{\text{max}} = 29,37 \times 3,6 = \mathbf{106\text{ km/h}} \]
\textbf{Resposta:} A velocidade máxima permitida é $\mathbf{29,4\text{ m/s}}$ (ou $\mathbf{106\text{ km/h}}$).

%------------------------------------------------------------------------------
\vspace{0.3cm}
\subsection*{Questão 25 [Sears 5.47]}
\textbf{Enunciado:} \textbf{Pêndulo Cônico.} Uma pequena esfera pesada de massa $m$ está presa à extremidade de uma corda leve de comprimento $L = 1,20\text{ m}$. A esfera gira com velocidade constante em uma trajetória circular horizontal de raio $r = 0,50\text{ m}$, de modo que a corda descreve a superfície de um cone reto formando um ângulo $\beta$ com a vertical. a) Determine o valor do ângulo $\beta$. b) Calcule a velocidade escalar constante da esfera ao longo do círculo. c) Determine o período de revolução $T$ do movimento circular.

\begin{dclbox}
\textbf{Geometria:} Triângulo retângulo formado pelo fio $L$, raio $r$ e altura vertical $h$:\\
$\sin\beta = \frac{r}{L}$.\\
\textbf{Forças sobre a esfera:}\\
1. Tensão $\vec{T}$ da corda (ao longo do fio, formando $\beta$ com a vertical).\\
   $T_y = T\cos\beta$ (vertical para cima); $T_{\text{rad}} = T\sin\beta$ (horizontal para o centro).\\
2. Peso $\vec{P} = mg$ (vertical para baixo).
\end{dclbox}

\noindent
\textbf{Resolução:}\\
a) Geometria do pêndulo:
\[ \sin\beta = \frac{r}{L} = \frac{0,50\text{ m}}{1,20\text{ m}} = 0,4167 \implies \beta = \arcsin(0,4167) = \mathbf{24,6^\circ} \]
b) No eixo vertical (não há aceleração vertical):
\[ T \cos\beta = mg \implies T = \frac{mg}{\cos\beta} \]
No eixo horizontal (radial para o centro da órbita):
\[ T \sin\beta = m \frac{v^2}{r} \]
Dividindo a equação radial pela vertical:
\[ \frac{T\sin\beta}{T\cos\beta} = \frac{m v^2 / r}{mg} \implies \tan\beta = \frac{v^2}{rg} \implies v = \sqrt{rg \tan\beta} \]
\[ v = \sqrt{(0,50\text{ m})(9,80\text{ m/s}^2)\tan(24,6^\circ)} = \sqrt{4,90 \times 0,4583} = \sqrt{2,246} = \mathbf{1,50\text{ m/s}} \]
c) O período de revolução $T_{\text{rev}}$ é o tempo para percorrer o perímetro $2\pi r$:
\[ T_{\text{rev}} = \frac{2\pi r}{v} = \frac{2\pi (0,50\text{ m})}{1,50\text{ m/s}} = \frac{\pi}{1,50} = \mathbf{2,09\text{ s}} \]
\textbf{Resposta:} a) $\beta = \mathbf{24,6^\circ}$. b) $v = \mathbf{1,50\text{ m/s}}$. c) $T_{\text{rev}} = \mathbf{2,09\text{ s}}$.

\newpage
%%%%%%%%%%%%%%%%%%%%%%%%%%%%%%%%%%%%%%%%%%%%%%%%%%%%%%%%%%%%%%%%%%%%%%%%%%%%%%%
\section*{PARTE III: RESOLUÇÃO DAS QUESTÕES DIFÍCEIS (26 A 40)}
%%%%%%%%%%%%%%%%%%%%%%%%%%%%%%%%%%%%%%%%%%%%%%%%%%%%%%%%%%%%%%%%%%%%%%%%%%%%%%%

%------------------------------------------------------------------------------
\subsection*{Questão 26 [Sears 4.45]}
\textbf{Enunciado:} Para estudar o dano que a colisão com grandes pássaros pode causar a um avião, projeta-se uma arma de teste pneumática que acelera projéteis ao longo de um cano cilíndrico reto, de modo que a posição do projétil em função do tempo é dada por:
\[ x(t) = (9,0 \times 10^3\text{ m/s}^2)\,t^2 - (8,0 \times 10^4\text{ m/s}^3)\,t^3 \]
O projétil de massa $m = 1,50\text{ kg}$ deixa a boca do cano no instante $t = 0,025\text{ s}$. a) Qual deve ser o comprimento do cano da arma de teste? b) Qual é a velocidade escalar do projétil no instante em que ele emerge do cano? c) Qual é o módulo da força resultante que atua sobre o projétil em: i) $t = 0$; ii) $t = 0,025\text{ s}$?

\noindent
\textbf{Resolução:}\\
a) O comprimento do cano é a posição do projétil no instante de saída $t = 0,025\text{ s}$:
\[ x(0,025) = (9,0 \times 10^3)(0,025)^2 - (8,0 \times 10^4)(0,025)^3 \]
\[ x(0,025) = (9000)(6,25 \times 10^{-4}) - (80000)(1,5625 \times 10^{-5}) = 5,625\text{ m} - 1,250\text{ m} = \mathbf{4,38\text{ m}} \]
b) A velocidade é obtida derivando a posição no tempo:
\[ v(t) = \frac{dx}{dt} = 2(9,0 \times 10^3)\,t - 3(8,0 \times 10^4)\,t^2 = 18\,000\,t - 240\,000\,t^2 \]
Em $t = 0,025\text{ s}$:
\[ v(0,025) = 18\,000(0,025) - 240\,000(0,025)^2 = 450 - 240\,000(6,25 \times 10^{-4}) = 450 - 150 = \mathbf{300\text{ m/s}} \]
c) A aceleração é obtida derivando a velocidade:
\[ a(t) = \frac{dv}{dt} = 18\,000 - 480\,000\,t \]
Pela Segunda Lei de Newton: $F_{\text{res}}(t) = m a(t) = 1,50 \times (18\,000 - 480\,000\,t)$:\\
i) Para $t = 0$:
\[ F_{\text{res}}(0) = 1,50\text{ kg} \times 18\,000\text{ m/s}^2 = \mathbf{2,70 \times 10^4\text{ N}} \]
ii) Para $t = 0,025\text{ s}$:
\[ a(0,025) = 18\,000 - 480\,000(0,025) = 18\,000 - 12\,000 = 6000\text{ m/s}^2 \]
\[ F_{\text{res}}(0,025) = 1,50\text{ kg} \times 6000\text{ m/s}^2 = \mathbf{9,00 \times 10^3\text{ N}} \]
\textbf{Resposta:} a) $L = \mathbf{4,38\text{ m}}$. b) $v = \mathbf{300\text{ m/s}}$. c) i) $\mathbf{2,70 \times 10^4\text{ N}}$; ii) $\mathbf{9,00 \times 10^3\text{ N}}$.

%------------------------------------------------------------------------------
\vspace{0.3cm}
\subsection*{Questão 27 [Sears 4.47]}
\textbf{Enunciado:} Um instrumento científico de $6,50\text{ kg}$ está pendurado por um cabo vertical no interior de uma espaçonave de teste que está sendo lançada verticalmente a partir da superfície terrestre. O foguete parte do repouso e atinge a altitude de $276\text{ m}$ em $15,0\text{ s}$ com aceleração vertical constante. a) Desenhe o diagrama do corpo livre para o instrumento durante essa fase do voo, indicando claramente qual força possui maior módulo. b) Calcule o módulo da força de tração que o cabo vertical exerce sobre o instrumento.

\begin{dclbox}
\textbf{Corpo:} Instrumento ($m = 6,50\text{ kg}$).\\
Forças verticais:\\
1. Tração $\vec{T}$: exercida pelo cabo, para cima.\\
2. Peso $\vec{P} = mg$: atração gravitacional terrestre, para baixo.\\
Como o foguete acelera para cima, a força resultante deve ser para cima: $\mathbf{T > P}$.
\end{dclbox}

\noindent
\textbf{Resolução:}\\
a) Pela cinemática do lançamento vertical ($v_0 = 0$, $y = 276\text{ m}$, $t = 15,0\text{ s}$):
\[ y = \frac{1}{2} a t^2 \implies a = \frac{2y}{t^2} = \frac{2(276\text{ m})}{(15,0\text{ s})^2} = \frac{552}{225} = 2,453\text{ m/s}^2 \]
b) Pela Segunda Lei de Newton aplicada no instrumento ($+y$ para cima):
\[ \sum F_y = T - mg = m a \implies T = m(g + a) \]
\[ T = 6,50\text{ kg} \times (9,80 + 2,453)\text{ m/s}^2 = 6,50 \times 12,253 = \mathbf{79,6\text{ N}} \]
\textbf{Resposta:} a) $T > P$ (a tração para cima supera o peso de $63,7\text{ N}$). b) A tração é $\mathbf{79,6\text{ N}}$.

%------------------------------------------------------------------------------
\vspace{0.3cm}
\subsection*{Questão 28 [Sears 4.53]}
\textbf{Enunciado:} Um paraquedista com massa total de $75,0\text{ kg}$ pula de uma aeronave a grande altitude. Durante a queda vertical antes de abrir o paraquedas, a força de resistência do ar oposta à velocidade é dada por $F_{\text{res}} = -k v^2$, onde $k = 0,250\text{ kg/m}$. a) Obtenha o valor da velocidade terminal $v_t$ atingida pelo paraquedista. b) Determine a aceleração instantânea de queda quando a sua velocidade for igual a metade da velocidade terminal ($v = v_t / 2$).

\begin{dclbox}
\textbf{Corpo:} Paraquedista caindo verticalmente ($m = 75,0\text{ kg}$).\\
Eixo vertical apontando para baixo ($+y$ para baixo):\\
1. Peso $P = mg$ (para baixo).\\
2. Resistência do ar $F_{\text{ar}} = k v^2$ (para cima, oposta à velocidade).\\
Equação: $mg - k v^2 = m a \implies a = g - \frac{k}{m}v^2$.
\end{dclbox}

\noindent
\textbf{Resolução:}\\
a) A velocidade terminal $v_t$ ocorre quando a aceleração se anula ($a = 0$):
\[ mg - k v_t^2 = 0 \implies v_t = \sqrt{\frac{mg}{k}} \]
\[ v_t = \sqrt{\frac{75,0\text{ kg} \times 9,80\text{ m/s}^2}{0,250\text{ kg/m}}} = \sqrt{\frac{735}{0,250}} = \sqrt{2940} = \mathbf{54,2\text{ m/s}} \quad (\approx 195\text{ km/h}) \]
b) No instante em que $v = \frac{v_t}{2}$:
\[ F_{\text{ar}} = k \left(\frac{v_t}{2}\right)^2 = \frac{1}{4} (k v_t^2) = \frac{1}{4} mg \]
Segunda Lei de Newton:
\[ m a = mg - F_{\text{ar}} = mg - \frac{1}{4} mg = \frac{3}{4} mg \implies a = \frac{3}{4} g = \frac{3}{4}(9,80\text{ m/s}^2) = \mathbf{7,35\text{ m/s}^2} \]
\textbf{Resposta:} a) $v_t = \mathbf{54,2\text{ m/s}}$. b) $a = \mathbf{7,35\text{ m/s}^2}$ para baixo.

%------------------------------------------------------------------------------
\vspace{0.3cm}
\subsection*{Questão 29 [Sears 4.59]}
\textbf{Enunciado:} Um objeto com massa $m$ move-se em linha reta ao longo do eixo $Ox$. Sua coordenada de posição em função do tempo é descrita por $x(t) = A t - B t^3$, onde $A$ e $B$ são constantes positivas conhecidas. a) Obtenha as expressões algébricas para a velocidade $v(t)$ e a aceleração $a(t)$ do objeto. b) Encontre a expressão analítica da força resultante $F(t)$ exercida sobre o corpo. c) Em qual instante de tempo a força resultante sobre o corpo se anula?

\noindent
\textbf{Resolução:}\\
a) Derivando sucessivamente a função posição:
\[ v(t) = \frac{dx}{dt} = \mathbf{A - 3Bt^2} \]
\[ a(t) = \frac{dv}{dt} = \mathbf{-6Bt} \]
b) Pela Segunda Lei de Newton:
\[ F_{\text{res}}(t) = m a(t) = \mathbf{-6mBt} \]
c) Para que a força resultante seja nula:
\[ F_{\text{res}}(t) = 0 \implies -6mBt = 0 \implies \mathbf{t = 0} \]
\textbf{Resposta:} a) $v(t) = A - 3Bt^2$; $a(t) = -6Bt$. b) $F(t) = -6mBt$. c) Em $t = \mathbf{0}$.

%------------------------------------------------------------------------------
\vspace{0.3cm}
\subsection*{Questão 30 [Sears 5.57]}
\textbf{Enunciado:} Um trabalhador de $80,0\text{ kg}$ está de pé sobre uma plataforma suspensa de $25,0\text{ kg}$ e puxa uma corda leve que passa por uma polia presa ao teto, de modo que ele e a plataforma aceleram juntos para cima com taxa constante de $1,20\text{ m/s}^2$. Despreze o atrito e as massas das polias e da corda. a) Desenhe os diagramas do corpo livre separados para o homem e para a plataforma. b) Calcule a força de tração que o homem deve exercer sobre a corda. c) Calcule o módulo da força normal que atua entre os pés do homem e a plataforma.

\begin{dclbox}
\textbf{Análise do Sistema Homem + Plataforma:}\\
Massa total: $M = m_h + m_p = 80,0 + 25,0 = 105,0\text{ kg}$.\\
A corda passa pela polia no teto e desce: um lado está preso à plataforma e o outro é puxado pelas mãos do homem. Portanto, \textbf{duas seções de corda sustentam o sistema para cima}, cada uma com tensão $T$.\\
\textbf{DCL do Homem isolado:} Tração $T$ (para cima, nas mãos), Normal $n$ (para cima, nos pés), Peso $m_h g$ (para baixo).\\
\textbf{DCL da Plataforma isolada:} Tração $T$ (para cima, no gancho), Normal $n$ (para baixo, exercida pelos pés do homem), Peso $m_p g$ (para baixo).
\end{dclbox}

\noindent
\textbf{Resolução:}\\
b) Aplicando a Segunda Lei ao conjunto $(homem + plataforma)$ acelerando para cima com $a = 1,20\text{ m/s}^2$:
\[ \sum F_y = 2T - (m_h + m_p)g = (m_h + m_p)a \]
\[ 2T = (m_h + m_p)(g + a) = 105,0\text{ kg} \times (9,80 + 1,20)\text{ m/s}^2 = 105,0 \times 11,0 = 1155\text{ N} \]
\[ T = \frac{1155\text{ N}}{2} = \mathbf{578\text{ N}} \]
c) Aplicando a Segunda Lei na plataforma isolada ($m_p = 25,0\text{ kg}$, acelerando para cima):
\[ \sum F_y = T - n - m_p g = m_p a \implies n = T - m_p(g + a) \]
\[ n = 577,5\text{ N} - (25,0\text{ kg})(11,0\text{ m/s}^2) = 577,5 - 275 = \mathbf{303\text{ N}} \]
(Verificando no homem: $T + n - m_h g = 577,5 + 302,5 - 784 = 880 - 784 = 96\text{ N} = m_h a = 80 \times 1,20 = 96\text{ N}$).\\
\textbf{Resposta:} b) Tração: $\mathbf{578\text{ N}}$. c) Força normal nos pés: $\mathbf{303\text{ N}}$.

%------------------------------------------------------------------------------
\vspace{0.3cm}
\subsection*{Questão 31 [Sears 5.63]}
\textbf{Enunciado:} \textbf{Corda com massa.} Em problemas de mecânica, as massas dos cabos e cordas costumam ser desprezadas. Quando uma corda pesada e flexível de massa $M$ está suspensa horizontalmente entre dois apoios fixados à mesma altura, suas extremidades formam um ângulo $\theta$ com a linha horizontal. a) Determine a tensão $T_1$ nas extremidades da corda em função de $M$, $g$ e $\theta$. b) Determine a tensão $T_0$ no ponto mais baixo (ponto médio inferior) da corda. c) Explique por que é fisicamente impossível que o ângulo $\theta$ seja rigorosamente igual a zero para uma corda com massa sob a ação da gravidade.

\begin{dclbox}
\textbf{DCL da corda inteira:} Duas extremidades com trações simétricas $T_1$ a um ângulo $\theta$ acima da horizontal; Peso total $Mg$ para baixo.\\
\textbf{DCL da metade da corda:} Extremidade superior com $T_1$ a $\theta$; Ponto médio com tração puramente horizontal $T_0$; Peso da meia corda $\frac{M}{2}g$ para baixo.
\end{dclbox}

\noindent
\textbf{Resolução:}\\
a) Equilíbrio vertical da corda inteira:
\[ \sum F_y = 2 T_1 \sin\theta - Mg = 0 \implies \mathbf{T_1 = \frac{Mg}{2\sin\theta}} \]
b) Isolando a metade esquerda da corda:\\
Equilíbrio horizontal ($\sum F_x = 0$): a tensão horizontal $T_0$ no meio equilibra a componente horizontal de $T_1$:
\[ T_0 = T_1 \cos\theta = \left(\frac{Mg}{2\sin\theta}\right)\cos\theta = \mathbf{\frac{Mg}{2\tan\theta}} \]
c) Se o ângulo fosse rigorosamente nulo ($\theta = 0$):
\[ \sin(0) = 0 \implies T_1 = \frac{Mg}{0} \to \infty \]
Para sustentar o peso vertical de qualquer porção de corda, é indispensável haver uma componente vertical de força para cima ($T_1\sin\theta$). Se a corda ficasse perfeitamente reta na horizontal ($\theta=0$), nenhuma tração teria componente vertical para equilibrar a gravidade. Como qualquer corda real possui limite de ruptura finito, \textbf{é fisicamente impossível mantê-la sem curvatura (catenária)}.\\
\textbf{Resposta:} a) $T_1 = \frac{Mg}{2\sin\theta}$. b) $T_0 = \frac{Mg}{2\tan\theta}$. c) Exigiria tração infinita para sustentar o peso.

%------------------------------------------------------------------------------
\vspace{0.3cm}
\subsection*{Questão 32 [Sears 5.70]}
\textbf{Enunciado:} Um foguete de pesquisa de $25\,000\text{ kg}$ é lançado verticalmente da superfície terrestre com aceleração constante dirigida para cima de $2,50\text{ m/s}^2$. Durante a fase inicial do movimento, adote $g = 9,80\text{ m/s}^2$. a) Qual é o módulo da força total de propulsão exercida pelos motores sobre o foguete? b) Um sensor experimental de $15,0\text{ kg}$ está apoiado no assoalho da câmara de carga do foguete. Calcule a força normal que o piso exerce sobre o sensor durante essa subida acelerada.

\noindent
\textbf{Resolução:}\\
a) Pela Segunda Lei de Newton no foguete ($M = 25\,000\text{ kg}$):
\[ \sum F_y = F_{\text{prop}} - Mg = M a \implies F_{\text{prop}} = M(g + a) \]
\[ F_{\text{prop}} = 25\,000\text{ kg} \times (9,80 + 2,50)\text{ m/s}^2 = 25\,000 \times 12,30 = \mathbf{3,08 \times 10^5\text{ N}} \]
b) Aplicando a Segunda Lei ao sensor de $m = 15,0\text{ kg}$ no piso acelerado:
\[ n - mg = m a \implies n = m(g + a) = 15,0\text{ kg} \times 12,30\text{ m/s}^2 = \mathbf{185\text{ N}} \]
\textbf{Resposta:} a) $F_{\text{prop}} = \mathbf{3,08 \times 10^5\text{ N}}$. b) $n = \mathbf{185\text{ N}}$.

%------------------------------------------------------------------------------
\vspace{0.3cm}
\subsection*{Questão 33 [Sears 5.71]}
\textbf{Enunciado:} Dois blocos retangulares de massas $m_A = 2,0\text{ kg}$ e $m_B = 4,0\text{ kg}$ estão sobrepostos: o bloco A está apoiado diretamente sobre o bloco B. O coeficiente de atrito estático entre as superfícies em contato dos blocos A e B é $\mu_e = 0,30$. A superfície da mesa horizontal sobre a qual o bloco B se apoia é perfeitamente lisa (sem atrito). Uma força horizontal $F$ é aplicada diretamente sobre o bloco inferior B. a) Qual é a aceleração máxima que o bloco A pode ter antes de começar a escorregar sobre o bloco B? b) Qual é o módulo máximo da força horizontal $F$ que pode ser aplicada ao bloco B para que os blocos A e B se acelerem juntos sem que haja deslizamento entre eles?

\begin{dclbox}
\textbf{DCL do bloco A (superior, $m_A = 2,0\text{ kg}$):}\\
Vertical: Peso $P_A = m_A g$ e Normal $n_A = m_A g$.\\
Horizontal: A \textbf{única força horizontal} que atua sobre A é a força de atrito estático exercida pelo bloco B: $f_{e,A}$ para a frente.\\
\textbf{DCL do bloco B (inferior, $m_B = 4,0\text{ kg}$):}\\
Força externa $F$ para a direita; força de atrito de reação de A sobre B ($f_{e,A}$ para trás).
\end{dclbox}

\noindent
\textbf{Resolução:}\\
a) A força horizontal máxima que o bloco B pode exercer sobre A é o atrito estático máximo:
\[ f_{e,\text{max}} = \mu_e n_A = \mu_e m_A g \]
Pela 2ª Lei de Newton sobre o bloco A:
\[ f_{e,\text{max}} = m_A a_{\text{max}} \implies a_{\text{max}} = \frac{\mu_e m_A g}{m_A} = \mu_e g = 0,30 \times 9,80\text{ m/s}^2 = \mathbf{2,94\text{ m/s}^2} \]
b) Para acelerar os dois blocos juntos com a aceleração máxima $a_{\text{max}} = 2,94\text{ m/s}^2$ sem que haja deslizamento relativo:
\[ F_{\text{max}} = (m_A + m_B) a_{\text{max}} = (2,0\text{ kg} + 4,0\text{ kg}) \times 2,94\text{ m/s}^2 = 6,0 \times 2,94 = \mathbf{17,6\text{ N}} \]
\textbf{Resposta:} a) $a_{\text{max}} = \mathbf{2,94\text{ m/s}^2}$. b) $F_{\text{max}} = \mathbf{17,6\text{ N}}$.

%------------------------------------------------------------------------------
\vspace{0.3cm}
\subsection*{Questão 34 [Sears 5.77]}
\textbf{Enunciado:} Um bloco com massa $m = 3,00\text{ kg}$ está sobre uma rampa inclinada de $37,0^\circ$ com a horizontal. O coeficiente de atrito cinético entre o bloco e a rampa é $\mu_c = 0,30$. O bloco é lançado para cima ao longo da rampa com velocidade inicial $v_0 = 8,00\text{ m/s}$. a) Calcule a distância que o bloco percorre rampa acima até parar momentaneamente. b) Verifique se o bloco, após parar no ponto mais alto, permanece em repouso ou desliza de volta para baixo, sabendo que $\mu_e = 0,40$. Em caso afirmativo, determine sua velocidade ao retornar à base.

\begin{dclbox}
\textbf{Fase 1: Subida do bloco}\\
Eixos: $+x$ apontando rampa acima. Forças no eixo $x$:\\
1. Componente paralela do peso: $P_x = -mg\sin(37,0^\circ)$ (puxa para baixo).\\
2. Força de atrito cinético: $f_c = \mu_c mg\cos(37,0^\circ)$ (oposta ao movimento, puxa para baixo!).\\
Ambas as forças somam-se desacelerando o bloco!
\end{dclbox}

\noindent
\textbf{Resolução:}\\
a) Aceleração durante a subida:
\[ \sum F_x = -mg\sin(37,0^\circ) - \mu_c mg\cos(37,0^\circ) = m a_{\text{subida}} \]
\[ a_{\text{subida}} = -g [\sin(37,0^\circ) + \mu_c \cos(37,0^\circ)] = -9,80 [0,6018 + 0,30(0,7986)] = -9,80(0,8414) = -8,246\text{ m/s}^2 \]
Pela Equação de Torricelli ($v = 0$):
\[ 0 = v_0^2 + 2 a d \implies d = \frac{v_0^2}{2 |a|} = \frac{(8,00)^2}{2(8,246)} = \frac{64,0}{16,492} = \mathbf{3,88\text{ m}} \]
b) No ponto mais alto ($v = 0$), testamos a tendência de escorregar rampa abaixo:\\
Tendência para baixo: $P_\parallel = mg\sin(37,0^\circ) = (3,00)(9,80)(0,6018) = 17,69\text{ N}$.\\
Atrito estático máximo: $f_{e,\text{max}} = \mu_e mg\cos(37,0^\circ) = 0,40(3,00)(9,80)(0,7986) = 9,39\text{ N}$.\\
Como $P_\parallel > f_{e,\text{max}}$ (isto é, $\tan 37^\circ \approx 0,75 > \mu_e = 0,40$), o atrito estático não consegue segurar o bloco, e ele \textbf{desliza de volta para baixo}.\\
Aceleração durante a descida (atrito aponta rampa acima):
\[ a_{\text{descida}} = g [\sin(37,0^\circ) - \mu_c \cos(37,0^\circ)] = 9,80 [0,6018 - 0,2396] = 9,80(0,3622) = 3,55\text{ m/s}^2 \]
Velocidade ao retornar à base ($v_0 = 0$, $d = 3,88\text{ m}$):
\[ v_{\text{base}} = \sqrt{2 a_{\text{descida}} d} = \sqrt{2 \times 3,55 \times 3,88} = \sqrt{27,55} = \mathbf{5,25\text{ m/s}} \]
\textbf{Resposta:} a) $d = \mathbf{3,88\text{ m}}$. b) O bloco desliza de volta e atinge a base com $v = \mathbf{5,25\text{ m/s}}$.

%------------------------------------------------------------------------------
\vspace{0.3cm}
\subsection*{Questão 35 [Sears 5.85]}
\textbf{Enunciado:} Três blocos estão conectados por meio de fios leves e polias ideais sem atrito: o bloco A de massa $m_A = 4,0\text{ kg}$ está suspenso verticalmente na extremidade esquerda; o bloco B de massa $m_B = 12,0\text{ kg}$ desliza sobre uma mesa horizontal com coeficiente de atrito cinético $\mu_c = 0,25$; e o bloco C está suspenso verticalmente na extremidade direita. O bloco B está se movendo para a direita com uma aceleração constante de $2,00\text{ m/s}^2$. a) Determine a tensão na corda que conecta o bloco A ao bloco B. b) Determine a massa necessária para o bloco pendente C.

\begin{dclbox}
\textbf{Movimento do sistema:} Bloco A sobe ($+a$); Bloco B acelera para a direita ($+a$); Bloco C desce ($+a$).\\
\textbf{Bloco A ($4,0\text{ kg}$):} $T_A - m_A g = m_A a$.\\
\textbf{Bloco B ($12,0\text{ kg}$):} $T_C - T_A - f_c = m_B a$, onde $f_c = \mu_c m_B g$.\\
\textbf{Bloco C ($m_C$):} $m_C g - T_C = m_C a$.
\end{dclbox}

\noindent
\textbf{Resolução:}\\
a) Pela equação do bloco A (que está subindo com $a = 2,00\text{ m/s}^2$):
\[ T_A = m_A(g + a) = 4,0\text{ kg} \times (9,80 + 2,00)\text{ m/s}^2 = 4,0 \times 11,80 = \mathbf{47,2\text{ N}} \]
b) Força de atrito sobre o bloco B na mesa horizontal:
\[ f_c = \mu_c m_B g = 0,25 \times (12,0\text{ kg} \times 9,80\text{ m/s}^2) = 29,4\text{ N} \]
Pela equação do bloco B:
\[ T_C = T_A + f_c + m_B a = 47,2\text{ N} + 29,4\text{ N} + (12,0\text{ kg})(2,00\text{ m/s}^2) = 76,6 + 24,0 = 100,6\text{ N} \]
Pela equação do bloco C (que está descendo acelerado):
\[ m_C g - T_C = m_C a \implies m_C (g - a) = T_C \]
\[ m_C = \frac{T_C}{g - a} = \frac{100,6\text{ N}}{9,80 - 2,00\text{ m/s}^2} = \frac{100,6}{7,80} = \mathbf{12,9\text{ kg}} \]
\textbf{Resposta:} a) $T_A = \mathbf{47,2\text{ N}}$. b) Massa do bloco C: $\mathbf{12,9\text{ kg}}$.

%------------------------------------------------------------------------------
\vspace{0.3cm}
\subsection*{Questão 36 [Sears 5.89]}
\textbf{Enunciado:} Dois objetos com massas de $5,0\text{ kg}$ e $2,0\text{ kg}$ estão suspensos a $0,600\text{ m}$ acima do solo, presos nas extremidades opostas de uma corda leve de $6,0\text{ m}$ que passa sobre uma polia fixa sem atrito. Os dois corpos partem do repouso. O objeto mais pesado desce até colidir com o solo e permanecer em repouso. Calcule a altura máxima que o objeto de $2,0\text{ kg}$ atinge acima do solo.

\noindent
\textbf{Resolução:}\\
\textbf{Fase 1 (Movimento acoplado enquanto o bloco pesado desce):}\\
$m_1 = 5,0\text{ kg}$ desce e $m_2 = 2,0\text{ kg}$ sobe. A aceleração conjunta é:
\[ a = \frac{m_1 - m_2}{m_1 + m_2} g = \frac{5,0 - 2,0}{5,0 + 2,0} \times 9,80 = \frac{3,0}{7,0} \times 9,80 = 4,20\text{ m/s}^2 \]
Quando o bloco 1 desce $h = 0,600\text{ m}$ e atinge o chão:
\begin{itemize}
    \item A altura do bloco 2 em relação ao solo é: $y_1 = 0,600\text{ m} + 0,600\text{ m} = 1,200\text{ m}$.
    \item A velocidade que o bloco 2 adquiriu é dada por:
    \[ v^2 = 2 a h = 2(4,20\text{ m/s}^2)(0,600\text{ m}) = 5,04\text{ m}^2/\text{s}^2 \implies v = 2,245\text{ m/s} \]
\end{itemize}
\textbf{Fase 2 (Voo livre do bloco 2 após a colisão do bloco 1):}\\
Com o bloco 1 em repouso no chão, a corda afrouxa ($T = 0$). O bloco 2 sobe em queda livre sob desaceleração $g = 9,80\text{ m/s}^2$:
\[ \Delta h_{\text{extra}} = \frac{v^2}{2g} = \frac{5,04}{2(9,80)} = \frac{5,04}{19,6} = 0,257\text{ m} \]
Altura máxima atingida pelo bloco de $2,0\text{ kg}$ em relação ao solo:
\[ H_{\text{max}} = y_1 + \Delta h_{\text{extra}} = 1,200\text{ m} + 0,257\text{ m} = \mathbf{1,46\text{ m}} \]
\textbf{Resposta:} A altura máxima é de $\mathbf{1,46\text{ m}}$ acima do solo.

%------------------------------------------------------------------------------
\vspace{0.3cm}
\subsection*{Questão 37 [Sears 5.91]}
\textbf{Enunciado:} Um bloco A com massa $m$ é pressionado contra a face vertical dianteira de um carrinho de transporte. O coeficiente de atrito estático entre o bloco e a face vertical do carrinho é $\mu_e$. a) Qual deve ser a aceleração mínima do carrinho para que o bloco A não escorregue para baixo em relação à parede vertical do carrinho? b) Descreva como as forças sobre o bloco são justificadas e compreendidas por um passageiro situado em repouso dentro do referencial do carrinho acelerado.

\begin{dclbox}
\textbf{Corpo:} Bloco A (em contato com a parede vertical).\\
1. Na horizontal: a parede vertical empurra o bloco para a frente com a força normal $n$.\\
2. Na vertical: a gravidade puxa o bloco para baixo com o peso $mg$; a força de atrito estático $f_e$ aponta para cima impedindo o deslizamento.
\end{dclbox}

\noindent
\textbf{Resolução:}\\
a) No referencial inercial da pista:
\[ \sum F_x = n = m a \]
No eixo vertical, para que o bloco permaneça em equilíbrio:
\[ \sum F_y = f_e - mg = 0 \implies f_e = mg \]
Como a força de atrito estático possui um limite máximo $f_e \le \mu_e n$:
\[ mg \le \mu_e n = \mu_e (m a) \implies g \le \mu_e a \implies a \ge \frac{g}{\mu_e} \]
\[ a_{\text{min}} = \mathbf{\frac{g}{\mu_e}} \]
b) No referencial do carrinho (não inercial com aceleração $+a\hat{\imath}$):\\
O passageiro observa o bloco parado em repouso relativo na parede. Para aplicar a 1ª Lei de Newton nesse referencial acelerado, ele introduz uma \textbf{força inercial fictícia} dirigida para trás: $\vec{F}_{\text{inercial}} = -m\vec{a}$. Essa força inercial empurra o bloco contra a parede, sendo equilibrada pela força normal da parede ($n = ma$). A força normal resultante gera a capacidade máxima de atrito $\mu_e n$ que equilibra exatamente o peso vertical $mg$.\\
\textbf{Resposta:} a) $a_{\text{min}} = \mathbf{\frac{g}{\mu_e}}$. b) No referencial acelerado, a força inercial $-m\vec{a}$ pressiona o bloco contra a parede gerando a força normal de suporte.

%------------------------------------------------------------------------------
\vspace{0.3cm}
\subsection*{Questão 38 [Sears 5.95]}
\textbf{Enunciado:} \textbf{Curva inclinada com atrito.} Uma curva com raio de curvatura igual a $120\text{ m}$ em uma estrada possui uma sobrelevação (ângulo de inclinação lateral $\theta$) projetada para que veículos a uma velocidade escalar de $20,0\text{ m/s}$ contornem a curva sem depender de nenhuma força de atrito lateral. a) Calcule o ângulo de sobrelevação lateral $\theta$ da rodovia. b) Se um veículo contornar essa mesma curva com velocidade de $30,0\text{ m/s}$, qual deve ser o coeficiente de atrito estático mínimo ($\mu_e$) entre os pneus e a pista para que o carro não derrape para fora da curva?

\begin{dclbox}
\textbf{Eixos adequados para curva inclinada:} Eixo horizontal radial (apontando para o centro da curva) e eixo vertical.\\
Força Normal $\vec{n}$: inclinada de $\theta$ com a vertical ($n_y = n\cos\theta$, $n_{\text{rad}} = n\sin\theta$).\\
Para velocidade superior ($v = 30\text{ m/s}$), o carro tende a subir a rampa $\implies$ Força de atrito estático aponta \textbf{descendo a rampa} (inclinada de $\theta$ abaixo da horizontal):\\
$f_{\text{rad}} = f_e \cos\theta$; $f_y = -f_e \sin\theta$.
\end{dclbox}

\noindent
\textbf{Resolução:}\\
a) Sem atrito a $v_1 = 20,0\text{ m/s}$:
\[ n\cos\theta = mg \quad \text{e} \quad n\sin\theta = m \frac{v_1^2}{R} \implies \tan\theta = \frac{v_1^2}{Rg} \]
\[ \tan\theta = \frac{(20,0)^2}{120 \times 9,80} = \frac{400}{1176} = 0,3401 \implies \theta = \mathbf{18,8^\circ} \]
b) A $v_2 = 30,0\text{ m/s}$ no limiar de escorregamento ($f_e = \mu_e n$):\\
Eixo vertical:
\[ n\cos\theta - f_e\sin\theta = mg \implies n(\cos\theta - \mu_e\sin\theta) = mg \]
Eixo radial horizontal:
\[ n\sin\theta + f_e\cos\theta = m \frac{v_2^2}{R} \implies n(\sin\theta + \mu_e\cos\theta) = m \frac{v_2^2}{R} \]
Dividindo as equações e simplificando por $\cos\theta$:
\[ \frac{\tan\theta + \mu_e}{1 - \mu_e\tan\theta} = \frac{v_2^2}{Rg} = \frac{(30,0)^2}{120 \times 9,80} = \frac{900}{1176} = 0,7653 \]
Substituindo $\tan\theta = 0,3401$:
\[ 0,3401 + \mu_e = 0,7653(1 - 0,3401\,\mu_e) = 0,7653 - 0,2603\,\mu_e \]
\[ \mu_e(1 + 0,2603) = 0,7653 - 0,3401 \implies 1,2603\,\mu_e = 0,4252 \implies \mu_e = \frac{0,4252}{1,2603} = \mathbf{0,337} \]
\textbf{Resposta:} a) $\theta = \mathbf{18,8^\circ}$. b) $\mu_e = \mathbf{0,337}$.

%------------------------------------------------------------------------------
\vspace{0.3cm}
\subsection*{Questão 39 [Sears 5.106]}
\textbf{Enunciado:} \textbf{Resistência de fluidos.} Uma pequena esfera pesada move-se no interior de um fluido viscoso, sofrendo uma força de resistência proporcional à velocidade escalar instantânea: $F_{\text{res}} = -k v$. Sabe-se que a velocidade terminal de descida da esfera é $v_t = 2,00\text{ m/s}$. A esfera é projetada verticalmente de baixo para cima dentro do fluido a partir da origem com velocidade inicial $v_0 = 6,00\text{ m/s}$. Despreze o empuxo. a) Na ausência de resistência do fluido, qual seria a altura máxima atingida pela esfera e quanto tempo ela levaria para atingir o topo? b) Levando em conta a resistência proporcional do fluido, deduza a equação da velocidade em função do tempo $v(t)$ e calcule a altura máxima atingida pela esfera no fluido.

\noindent
\textbf{Resolução:}\\
Velocidade terminal na descida: $mg = k v_t \implies \frac{k}{m} = \frac{g}{v_t} = \frac{9,80}{2,00} = 4,90\text{ s}^{-1}$.\\
a) Na ausência de resistência do fluido:
\[ H = \frac{v_0^2}{2g} = \frac{(6,00)^2}{2(9,80)} = \frac{36,0}{19,6} = \mathbf{1,84\text{ m}} \]
\[ t_{\text{topo}} = \frac{v_0}{g} = \frac{6,00}{9,80} = \mathbf{0,612\text{ s}} \]
b) Com resistência do fluido subindo ($+y$ para cima, onde tanto a gravidade quanto o arrasto apontam para baixo):
\[ m \frac{dv}{dt} = -mg - kv \implies \frac{dv}{dt} = -g \left(1 + \frac{v}{v_t}\right) \]
Para achar a altura máxima $H$ diretamente em função da velocidade, usamos a regra da cadeia: $\frac{dv}{dt} = v \frac{dv}{dy}$:
\[ v \frac{dv}{dy} = -g \left(1 + \frac{v}{v_t}\right) \implies dy = -\frac{1}{g} \frac{v}{1 + v/v_t} dv = -\frac{v_t^2}{g} \frac{v/v_t}{1 + v/v_t} \frac{dv}{v_t} \]
Definindo $u = v/v_t$:
\[ dy = -\frac{v_t^2}{g} \left(1 - \frac{1}{1+u}\right) du \]
Integrando desde $y=0$ ($v=v_0$, $u_0 = v_0/v_t = 3$) até $y=H$ ($v=0$, $u=0$):
\[ H = \frac{v_t^2}{g} \int_0^{u_0} \frac{u}{1+u} du = \frac{v_t^2}{g} [ u_0 - \ln(1 + u_0) ] \]
Substituindo $u_0 = \frac{6,00}{2,00} = 3$ e $v_t = 2,00\text{ m/s}$:
\[ H = \frac{(2,00)^2}{9,80} [ 3 - \ln(1 + 3) ] = \frac{4,00}{9,80} [ 3 - \ln(4) ] = 0,4082 [ 3 - 1,3863 ] = 0,4082 \times 1,6137 = \mathbf{0,659\text{ m}} \]
\textbf{Resposta:} a) Sem resistência: $H = \mathbf{1,84\text{ m}}$ e $t = \mathbf{0,612\text{ s}}$. b) Com resistência: $H = \mathbf{0,659\text{ m}}$ (redução expressiva de mais de 64\% na altura devido ao amortecimento viscoso).

%------------------------------------------------------------------------------
\vspace{0.3cm}
\subsection*{Questão 40 [Sears 5.121]}
\textbf{Enunciado:} Uma cunha triangular de massa $M$ e ângulo de inclinação $\alpha$ repousa sobre uma superfície horizontal perfeitamente lisa sem atrito. Um pequeno bloco de massa $m$ é colocado sobre a rampa da cunha, também sem atrito entre o bloco e a cunha. Uma força horizontal constante $F$ é aplicada à cunha. Qual deve ser o módulo da força $F$ para que o bloco permaneça em repouso em relação à cunha (isto é, mantendo-se a uma altura constante em relação à base e sem deslizar rampa acima nem rampa abaixo)?

\begin{dclbox}
\textbf{Corpo:} Bloco de massa $m$.\\
O bloco deve mover-se solidário com a cunha, possuindo aceleração puramente horizontal $\vec{a} = a\hat{\imath}$.\\
Forças sobre o bloco:\\
1. Força Peso $\vec{P} = mg$ (vertical para baixo).\\
2. Força Normal $\vec{n}$ da rampa: perpendicular à superfície da cunha.\\
Como a rampa faz ângulo $\alpha$ com a horizontal, a reta normal à rampa forma ângulo $\alpha$ com a \textbf{vertical}.
\end{dclbox}

\noindent
\textbf{Resolução:}\\
Decompondo as forças que atuam sobre o bloco $m$ em eixos horizontal e vertical:\\
Eixo vertical (o bloco mantém altura constante, $a_y = 0$):
\[ \sum F_y = n \cos\alpha - mg = 0 \implies n = \frac{mg}{\cos\alpha} \]
Eixo horizontal (o bloco acelera para a frente com a mesma aceleração $a$ da cunha):
\[ \sum F_x = n \sin\alpha = m a \]
Substituindo $n = \frac{mg}{\cos\alpha}$:
\[ \left(\frac{mg}{\cos\alpha}\right)\sin\alpha = m a \implies mg \tan\alpha = m a \implies a = g \tan\alpha \]
Agora, aplicando a Segunda Lei de Newton ao sistema completo formado pela cunha e pelo bloco de massa total $(M + m)$ sobre a mesa horizontal sem atrito:
\[ F = (M + m) a \]
Substituindo a aceleração horizontal necessária $a = g\tan\alpha$:
\[ \mathbf{F = (M + m) g \tan\alpha} \]

\begin{dica}
\textbf{Análise física dos limites:}
\begin{itemize}
    \item Se $\alpha \to 0^\circ$ (chão horizontal): $\tan 0^\circ = 0 \implies F = 0$ (o bloco fica em repouso sem precisar de força acelerando a cunha).
    \item Se $\alpha \to 90^\circ$ (parede vertical): $\tan 90^\circ \to \infty$ (é impossível sustentar um bloco sem atrito em uma parede puramente vertical apenas empurrando horizontalmente).
\end{itemize}
\end{dica}

\noindent
\textbf{Resposta:} A força horizontal deve ter módulo $\mathbf{F = (M + m) g \tan\alpha}$.

\end{document}
